\section*{第 \S\arabic{exercisegroup} 节习题}
\phantomsection
\addcontentsline{toc}{section}{第 \protect\S\arabic{exercisegroup} 节习题}

\(\mathfrak{J}1)\) a) 设 \(g\) 是 \(x \geqslant 0\) 上的连续实值函数。证明：若 \(g\) 满足两个恒等式

\[
\sum_ {k = 0} ^ {p - 1} g \left(\frac {x + k}{p}\right) = g (x)
\]

\[
\sum_ {h = 0} ^ {q - 1} g \left(\frac {x + h}{q}\right) = g (x)
\]

则它也满足恒等式

\[
\sum_ {j = 0} ^ {p q - 1} g \left(\frac {x + j}{p q}\right) = g (x).
\]

\begin{enumerate}
\def\labelenumi{\alph{enumi})}
\setcounter{enumi}{1}
\tightlist
\item
  由此推出：若函数 \(g\) 在 \(x \geqslant 0\) 上有连续导数且满足恒等式
\end{enumerate}

\[
\sum_ {k = 0} ^ {p - 1} g \left(\frac {x + k}{p}\right) = g (x)
\]

则它具有 \(a(x-\frac{1}{2})\) 的形式，其中 a 是常数（注意 g 满足把 p 换成 \(p^{n}\) 后的类似恒等式；令 n 趋于 \(+\infty\) ，推出 \(g'(x)=\int_{0}^{1}g'(t)dt\) ）。c) 由 b) 断定：\(\Gamma\) 函数是在 x\textgreater0 上有连续导数、且对 p 的某个值满足 VII, p.~305 的方程 (1) 与 VII, p.~318 的倍乘公式 (12) 的唯一函数。

\begin{enumerate}
\def\labelenumi{\arabic{enumi})}
\setcounter{enumi}{1}
\tightlist
\item
  对每个整数 \(k > 1\) ，令 \(S_k = \sum_{n=1}^{\infty} n^{-k}\) 。证明：对 \(-1 < x \leqslant 1\) ，有
\end{enumerate}

\[
\log \Gamma (1 + x) = - \gamma x + \sum_ {k = 2} ^ {\infty} (- 1) ^ {k} \frac {\mathrm{S} _ {k}}{k} x ^ {k}
\]

右端的级数在 \(]-1,1]\) 的每个紧子集上一致收敛，且对 \(|x|<1\) 绝对收敛。

\begin{enumerate}
\def\labelenumi{\arabic{enumi})}
\setcounter{enumi}{2}
\tightlist
\item
  设 s 是固定的实数；证明：当 t 趋于 \(+\infty\) 或 \(-\infty\) 时，有
\end{enumerate}

\[
| \Gamma (s + i t) | \sim \sqrt {2 \pi} | t | ^ {s - 1 / 2} e ^ {- (\pi / 2) | t |}
\]

以及

\[
\frac {\Gamma^ {\prime} (s + i t)}{\Gamma (s + i t)} \sim \log | t |.
\]

\begin{enumerate}
\def\labelenumi{\arabic{enumi})}
\setcounter{enumi}{3}
\tightlist
\item
  设 t 是固定的实数 \(\neq 0\) ；证明：当 s 趋于 \(+\infty\) 时，
\end{enumerate}

\[
| \Gamma (- s + i t) | \sim \sqrt {\frac {\pi}{2}} \left(s ^ {s + 1 / 2} e ^ {- s} \sqrt {\sinh^ {2} \pi t + \sin^ {2} \pi s}\right) ^ {- 1}
\]

（利用余元公式）。

\begin{enumerate}
\def\labelenumi{\arabic{enumi})}
\setcounter{enumi}{4}
\tightlist
\item
  设 \(x_{n}\) 是方程 \(\Gamma'(x) = 0\) 在区间 \(] - n, -n + 1[\) 中的根。证明
\end{enumerate}

\[
x _ {n} = - n + \frac {1}{\log n} + O \left(\frac {1}{(\log n) ^ {2}}\right)
\]

（利用余元公式与 VII, p.~316 的公式 (5)）。由此推出

\[
\Gamma (x _ {n}) \sim \frac {(- 1) ^ {n}}{\sqrt {2 \pi}} n ^ {- n - 1 / 2} e ^ {n + 1} \log n.
\]

\begin{enumerate}
\def\labelenumi{\arabic{enumi})}
\setcounter{enumi}{5}
\tightlist
\item
  设 \(V_{n}\) 是范德蒙德行列式 \(\mathrm{V}(1,2,\dots,n)\) （Alg., III, p.~532）。证明
\end{enumerate}

\[
\log V _ {n} = \frac {n ^ {2}}{2} \log n - \frac {3 n ^ {2}}{4} + \left(\frac {1}{2} \log 2 \pi - \frac {1}{4}\right) n - \frac {1}{12} \log n + k + O \left(\frac {1}{n}\right)
\]

其中 k 是常数。

